% ARCHIVED AUDITED DRAFT: independently checked by root, 2026-11-19 15:46 UTC.
% Incorporated into part_B.tex, whose version is authoritative.
\begin{saresult}{B.063}{Proved}{Every valuation domain with arbitrary coefficient extensions}
Let $(V,\mathfrak m)$ be a nonfield valuation domain,
of arbitrary rank, and let $V\subset S$ be any
unital commutative extension. If the residue field
is infinite, every such extension has the unique
standard polynomial co-operations. If the residue
field is finite, the polynomial coaddition and
comultiplication tensors are uniquely determined
if and only if
\[
                         \operatorname{Ann}_S(\mathfrak m)=0.
                                                               \tag{B63.1}
\]
Every positive case has all multivariable tensor
isomorphisms. This removes the torsion-free
restriction from B.041; base torsion away from
the maximal-ideal socle is permitted.
\end{saresult}
\begin{proof}
Infinite residue field is the unit Vandermonde
case. Suppose the residue field is $k=\mathbb F_q$.
First let $\mathfrak m$ be nonprincipal. Its value
group has no least positive element. The elementary
ordered-group argument in B.041 gives, for every
$\gamma>0$ and integer $N$, a positive $\delta$
with $N\delta<\gamma$.
For a polynomial of degree at most $d\ge1$, choose
$a\in\mathfrak m$ of value $\delta$ and the grid
$0,a,a^2,\ldots,a^d$. Its Vandermonde determinant
has value $d(d+1)(d+2)\delta/6$. Cramer's rule
shows that this determinant sends every coefficient
of an integer-valued polynomial into $V$.
Taking $\delta$ small enough for each prescribed
$r\in\mathfrak m$ and using divisibility in a
valuation domain proves that every coefficient
belongs to
\[
                         C=(V:_S\mathfrak m).
\]
This is a ring because $\mathfrak m^2=\mathfrak m$:
factor an element of positive value into two
elements of smaller positive value.

Localizing $S$ at $V\setminus\{0\}$ maps each
$c\in C$ into the fraction field $K$, since
$rc\in V$ for any nonzero $r\in\mathfrak m$.
That rational image must lie in $V$: a negative
value is contradicted by multiplying by an element
of sufficiently small positive value. Subtracting
this element of $V$ leaves an element annihilated
by $\mathfrak m$, because its products with
$\mathfrak m$ are both in $V$ and zero after
localization. Thus, as $V$-modules,
\[
                    C=V\oplus T,\qquad
                    T=\operatorname{Ann}_S(\mathfrak m).
\]
Writing $Q(X)=X^q-X$, monic division over the
$k$-vector space $T$ gives
\[
                    \Omega(V,S)=V[X]\oplus Q(X)T[X].     \tag{B63.2}
\]
Indeed the second summand consists exactly of
the $T$-coefficient polynomials vanishing on $k$.
If $0\ne z\in T$, the tensor
\[
                  zQ\otimes Q-Q\otimes zQ
\]
has zero polynomial image and is nonzero: its
two terms lie in distinct direct summands of
the tensor square in (B63.2), and $Q$ is monic
in the free module $V[X]$. Hence uniqueness
fails. If $T=0$, the polynomial module is
$V[X]$, proving this case.

Now suppose $\mathfrak m=\pi V$. Its least positive
value is $v(\pi)$. Each quotient $V/\pi^nV$ has
$q^n$ elements. Choosing any finite grid with
distinct residues modulo a sufficiently high
$\pi$-power makes its Vandermonde determinant
a unit times a power of $\pi$. Thus every
integer-valued coefficient lies in
\[
             C=\{s\in S:\pi^ns\in V\text{ for some }n\}.
                                                               \tag{B63.3}
\]
If $\operatorname{Ann}_S(\pi)=0$, this ring embeds
in $V[1/\pi]\subset K$. Every element of
$V[1/\pi]\setminus V$ has value a negative
integer multiple of $v(\pi)$, so an overring
containing one such element contains $\pi^{-1}$.
Consequently $C=V$ or $C=V[1/\pi]$. The former
gives ordinary polynomials; B.041 gives all
tensor isomorphisms in the latter case.

For necessity, put
\[
                         P=\bigcap_{n\ge0}\pi^nV.
\]
An element outside $P$ is a unit times a
nonnegative power of $\pi$: its value is bounded
by a finite multiple of the least positive
value, so is exactly such an integer multiple.
It follows that $P$ is prime and $V/P$ is a
discrete valuation ring. The ideal $P$ remains
an ideal of $C$, and $PC=P$. To see this, if
$\pi^Nc=a\in V$ and $p\in P$, write
$p=\pi^{N+j}b_j$ for every $j$. Then
$pc=\pi^jb_ja\in\pi^jV$ for every $j$, as required.

Write $\bar V=V/P$, $\bar C=C/P$. Reduction of
coefficients gives a surjection
\[
 \Omega(V,C)\longrightarrow\Omega(\bar V,\bar C)
 \quad\text{with kernel }P[X]=P\Omega(V,C).              \tag{B63.4}
\]
Indeed lift the finitely many coefficients of
an integer-valued polynomial over $\bar C$.
Its values on $V$ lie in the preimage of
$\bar V$, which is $V$; hence the lift is
integer-valued. The kernel and last equality
follow from $PC=P$.

Tensor injectivity descends through (B63.4),
despite the possible nonflatness of $V/P$.
Given a tensor in the reduced kernel, lift it
to $z\in\Omega(V,C)^{\otimes2}$. Its polynomial
image belongs to $P[X,Y]$. That polynomial
has a tensor preimage formed from ordinary
monomials with coefficients in $P$, and this
preimage reduces to zero. Subtract it from $z$.
Injectivity upstairs forces the resulting tensor,
and thus the original reduced tensor, to vanish.
The two existence conditions descend as well,
since every reduced polynomial lifts.

Apply B.058 to the discrete valuation ring
$\bar V$; every element of $\bar C/\bar V$ is
killed by a power of $\pi$. It follows that
$\bar C=\bar V$ or $\bar C=\operatorname{Frac}(\bar V)$.
The first equality forces $C=V$. In the second
case lift an inverse of $\pi$: one has
$\pi c=1+p$ for some $p\in P$. Since $1+p$
is a unit of $V$, $\pi$ is a unit of $C$,
and (B63.3) gives $C=V[1/\pi]$.
Both possibilities have no $\pi$-torsion.
But every element of $\operatorname{Ann}_S(\pi)$
belongs to $C$, proving its vanishing.
Since $\operatorname{Ann}_S(\mathfrak m)=
\operatorname{Ann}_S(\pi)$, this proves necessity.

The multivariable coefficient assertion used in
the positive cases follows from rectangular
Vandermonde grids, or by substituting widely
separated powers of one variable to keep all
monomials distinct. Thus the positive tensor
isomorphisms compute the full multivariable
integer-valued modules, including the original
coefficient extension $S$.
\end{proof}

\begin{saresult}{B.064}{Proved}{A universal coefficient-ring obstruction without Noetherian hypotheses}
For any unital inclusion $R\subset S$ and $d\ge1$,
let $D_d(R)$ be the ideal generated by all
Vandermonde determinants on $d+1$ elements of $R$,
and put
\[
 C_d=(R:_S D_d(R)),\qquad C=\bigcup_{d\ge1}C_d.           \tag{B64.1}
\]
Then $C$ is a subring of $S$. It is exactly the
set of coefficients that occur in integer-valued
polynomials, and also exactly the set of their
possible leading coefficients, with zero included. Every multivariable
integer-valued polynomial has its coefficients in
$C$. If the second polynomial tensor map is
injective, the inclusion $R\to C$ is a ring
epimorphism. If $C=R$, all tensor maps are the
ordinary polynomial isomorphisms.
\end{saresult}
\begin{proof}
The ideals $D_d$ decrease with $d$, since a
larger Vandermonde product is a multiple of
one on a subset of its nodes. Thus the $C_d$
increase, and $C_1=R$ because the nodes $0,1$
have determinant one.
Cramer's rule says that every coefficient of
an integer-valued polynomial of degree at most
$d$ lies in $C_d$.

There is an explicit monic polynomial $P_d$
whose every value belongs to $D_d$. For $d=1$
take $P_1=1$. For $d\ge2$, take the Vandermonde
product on the polynomial nodes
\[
                     0,\ 1,\ X,\ X^2,\ldots,X^{d-1}.
\]
Every factor with positive degree is monic,
and the only degree-zero factor is $1$; hence
$P_d$ is monic over every base ring.
For any $c\in C_d$, the polynomial $cP_d$
takes $R$ into $R$ and has leading coefficient
$c$. This proves both coefficient characterizations.
Sums of leading coefficients are obtained after
multiplying the corresponding polynomials by
powers of $X$ to equalize their degrees; products
come from multiplication of polynomials.
Therefore their set $C$ is a subring.

For a multivariable polynomial, substitute
$X_i=T^{N_i}$ with sufficiently separated
positive exponents so that its finitely many
monomials remain distinct. The resulting
one-variable polynomial is integer-valued on
$R$, and has exactly the same coefficients.
Thus all coefficients lie in $C$.

For $c\in C_d$, both $P_d$ and $cP_d$ belong
to $\Omega(R,S)$. The tensor
$cP_d\otimes P_d-P_d\otimes cP_d$ has zero
polynomial image. If the tensor map is
injective it vanishes. Extract the top
coefficient in each factor, now as maps
$\Omega(R,S)\to C$, to get
$c\otimes1=1\otimes c$ in $C\otimes_R C$.
This for every $c\in C$ makes multiplication
$C\otimes_R C\to C$ an isomorphism.
Finally $C=R$ gives $\Omega_n=R[X_1,\ldots,X_n]$
by the coefficient statement, proving the
last assertion.
\end{proof}

\begin{saresult}{B.065}{Proved}{Every absolutely flat base with arbitrary coefficients}
Let $R$ be an absolutely flat commutative ring,
meaning that every $R$-module is flat, and let
$R\subset S$ be any unital commutative extension.
With the explicit Vandermonde ideals of B.064,
the following are equivalent:
\[
\begin{gathered}
 \text{the polynomial co-operations are uniquely determined};\\
 \theta_2\text{ is injective};\qquad
 (R:_S D_d(R))=R\text{ for every }d;\\
 \operatorname{Ann}_{S/R}(D_d(R))=0\text{ for every }d;\qquad
 \Omega(R,S)=R[X].                                       \tag{B65.1}
\end{gathered}
\]
Every positive case has all polynomial tensor
isomorphisms. In particular, for any Boolean
ring, finite or infinite, the positive
extensions are exactly $S=R$.
\end{saresult}
\begin{proof}
By B.064, injectivity makes $R\to C$ an
injective ring epimorphism. Such an epimorphism
over an absolutely flat base is an isomorphism.
Indeed $C/R$ is flat, so tensoring
$0\to R\to C\to C/R\to0$ with any module
$M$ keeps $M\to C\otimes_R M$ injective.
But the epimorphism condition gives
$C\otimes_R(C/R)=0$. Take $M=C/R$ to conclude
$C/R=0$.
Thus injectivity implies $C=R$. B.064 shows
that this is equivalent to the coefficient
conditions and to $\Omega=R[X]$, and gives
all positive tensor isomorphisms. The
annihilator condition is the same assertion
in the quotient module $S/R$.

For a Boolean ring, $x^2=x$ for every $x$.
Every prime quotient is the field with two
elements, so every three-node Vandermonde
determinant vanishes in every prime quotient.
The ring is reduced, hence $D_2(R)=0$.
Consequently $C_2=S$, and the criterion is
exactly $S=R$. Boolean rings are absolutely
flat: every finitely generated ideal is
generated by an idempotent, and the usual
finite-equation criterion for flatness then
shows that every module is flat. More
explicitly, for any finitely generated ideal
$I=eR$, the map $I\otimes_R M\to M$ identifies
with the inclusion $eM\hookrightarrow M$;
the ideal criterion proves flatness, as in the opened
Stacks Lemma 10.39.5, tag 00HD. The equivalent
reduced zero-dimensional characterization of absolute
flatness is Stacks Lemma 15.106.5, tag 092F
\cite{StacksAbsoluteFlat}.
\end{proof}
